\scanpage{130}
\uline{Corollaire}\quad\margin{\rule[-3.8cm]{0.4pt}{3.8cm}} Si $A\in\mathcal L^1(\mathcal H)$, soit
$M_A(\unsure{v|r})=\operatorname*{Sup}_{|z|=\unsure{v|r}}\det(1+zA)$
\deleted{si $A$ est \unsure{hermitien} \unsure{positif}, et si $A\in\mathcal L^p(\mathcal H)$}, \deleted{\unsure{tel|on}}
\deleted{\unsure{avec}} \inserted{et} $0<p<1$\,; on a
\[
S_p(A)=p\frac{\sin\pi p}{\pi}\int_0^\infty\frac{\log M_{|A|}(\unsure{v|r})}{{\unsure{v|r}}^{1+p}}\,d\unsure{v|r}\tag*{(9)}
\]

Il suffit en effet de poser $f(z)=\det(1+z|A|)$,
les $\alpha_i$ étant les valeurs propres de $|A|$\,; on aura
ici $M(\unsure{v|r})=f(\unsure{v|r})$.

Pour démontrer (4), on \deleted{\unsure{pe|peut}} peut écrire
\[
S_p(A+B)=S_p(|A+B|)=p\frac{\sin\pi p}{\pi}\int_0^\infty
\frac{\log M_{|A+B|}(\unsure{v|r})}{{\unsure{v|r}}^{1+p}}\,d\unsure{v|r}
\]
Or $M_{|A+B|}(\unsure{v|r})=\det(1+\unsure{v|r}|A+B|)\leq\det(1+\unsure{v|r}|A|)\det(1+\unsure{v|r}|B|)$
(formule \unsure{\ldots}) \deleted{$\leq\unsure{rM_A|M_A}$} soit
\[
M_{|A+B|}(\unsure{v|r})\leq M_{|A|}(\unsure{v|r})\cdot M_{|B|}(\unsure{v|r})\tag*{(10)}
\]
d'où
\[
S_p(A+B)\leq\frac{p\sin\pi p}{\pi}\int_0^\infty\frac{\log M_{|A|}(\unsure{v|r})}{{\unsure{v|r}}^{1+\unsure{v|r|p}}}\,d\unsure{v|r}
+\frac{p\sin\pi p}{\pi}\int_0^\infty\frac{\log M_{|B|}(\unsure{v|r})}{{\unsure{v|r}}^{1+\unsure{r|v|p}}}\,d\unsure{v|r}
\]
Appliquant de nouveau (9), on obtient bien la
formule voulue (4).

\deleted{\unsure{Pour prouver}}

2) Pour prouver la formule (5), je ne connais
qu'une méthode analytique. De façon générale,
\unsure{.} si $A$ est un opérateur cp\textsuperscript{t} continu, on prend la
suite des val. propres $(\alpha_i)$, \deleted{\unsure{prises|posons}} et si
$0<p<+\infty$, posons $s_p(A)=\sum|\alpha_i|^p$.
\deleted{\unsure{On a} $\unsure{s_p(A)|S_p(A)}\unsure{\leq|=}$ \unsure{$+\infty$|$S_p$} \unsure{évidemment pour $p=1$|évidemment pour tout entier} Si $n>0$ $s_p(A)=s_{\frac{p}{n}}(A^n)$. \unsure{Supposons}}
